\begin{proof}[Proof of \cref{obj:thm:supplied-propensity}]
Fix the public parameters in the stated domain. Throughout, inverse-propensity quotients are assigned value zero when their denominator is zero:
\[
\frac{t_{\mathrm{num}}}{0}:=0
\qquad(t_{\mathrm{num}}\in\mathbb R).
\]
The designated propensity is extended by zero outside \([0,1]^d\).

\begin{enumerate}
\item Define the constant-nuisance subfamily and its supplied-propensity risk by
\[
\mathcal M_{1/2}
=
\left\{
P\in\mathcal M:
e_P(x)=\mu_{0,P}(x)=\frac12
\text{ for every }x\in[0,1]^d
\right\},
\]
\[
\mathcal R^e_{n,m}(T^e,P)
=
\int
\left|T^e(\mathcal D_{n,m},U,e_P)-\tau_P(x_0)\right|
\,d\mathsf E_{n,m}(P),
\qquad
R^e_{1/2}(n,m)
=
\inf_{T^e}\sup_{P\in\mathcal M_{1/2}}
\mathcal R^e_{n,m}(T^e,P).
\]
These loss integrals may take the value \(+\infty\). Since the subfamily uses the same decision class as \cref{obj:def:oracle-risk}, restriction of the supremum gives
\[
R^e_{1/2}(n,m)\le R^e(n,m).
\]
It therefore suffices to establish a lower bound for \(R^e_{1/2}(n,m)\), a uniform upper bound for the displayed smoother, and its admissibility. We first construct the lower-bound family.
% lean: CausalSmith.Stat.TwosamplePointcateAnnotationFrontier.oracleSubfamilyRisk_le_oracleRisk

\item Use the smooth profile from \cref{obj:def:marked-handle}, with the explicit formulas
\[
\chi(\upsilon)=
\begin{cases}
\exp(-1/\upsilon),&\upsilon>0,\\
0,&\upsilon\le0,
\end{cases}
\qquad \upsilon\in\mathbb R,
\qquad
B(w)=\prod_{\iota_{\mathrm{cov}}=1}^d\frac{\chi(1-4w_{\iota_{\mathrm{cov}}}^2)}{\chi(1)},
\qquad w\in\mathbb R^d.
\]
For this argument, extend the scaled-bump formula to every positive localization width and every covariate vector:
\[
B_h(x)=B\bigl((x-x_0)/h\bigr),
\qquad
C_h=\prod_{j=1}^d
\left[\frac12-\frac h2,\frac12+\frac h2\right],
\qquad h>0,\quad x\in\mathbb R^d.
\]
Smoothness of the exponential glue function and finite products give
\[
B\in C^\infty(\mathbb R^d),
\qquad
\operatorname{supp}(B)\subseteq[-1/2,1/2]^d,
\qquad
0\le B\le1,
\qquad
B(0)=1.
\]
In particular,
\[
0\le B_h\le1,\qquad
B_h(x)=0\quad(x\notin C_h),\qquad
B_h(x_0)=1.
\]

We need a uniform Hölder bound after scaling. Define
\[
p_*=\lceil\gamma\rceil-1,
\qquad
s_o=\gamma-p_*,
\qquad
0\le p_*\le2,\quad 0<s_o\le1.
\]
Write the iterated derivatives as
\[
D^0B=B,\qquad D^{\iota_{\mathrm{der}}+1}B=D(D^{\iota_{\mathrm{der}}}B),
\qquad \iota_{\mathrm{der}}\in\mathbb N,
\]
and define their finite operator-norm bounds by
\[
C_{B,\iota_{\mathrm{der}}}=\sup_{w\in\mathbb R^d}\|D^{\iota_{\mathrm{der}}}B(w)\|,
\qquad \iota_{\mathrm{der}}\in\mathbb N.
\]
Their finiteness follows from continuity and compact support. Set
\[
A_B=\max\{C_{B,p_*},C_{B,p_*+1}\},
\qquad
M_B=2A_B.
\]
The derivative bound and the mean-value inequality give, for every \(w,w'\in\mathbb R^d\),
\[
\|D^{p_*}B(w)-D^{p_*}B(w')\|
\le A_B\|w-w'\|_2,
\qquad
\|D^{p_*}B(w)-D^{p_*}B(w')\|\le2A_B.
\]
If \(\|w-w'\|_2\le1\), use
\(\|w-w'\|_2\le\|w-w'\|_2^{s_o}\); if
\(\|w-w'\|_2>1\), use \(1\le\|w-w'\|_2^{s_o}\).
Thus both cases yield
\[
\|D^{p_*}B(w)-D^{p_*}B(w')\|
\le M_B\|w-w'\|_2^{s_o}.
\]

Translation and dilation satisfy
\[
D^{\iota_{\mathrm{der}}}B_h(x)=h^{-\iota_{\mathrm{der}}}D^{\iota_{\mathrm{der}}}B\bigl((x-x_0)/h\bigr).
\]
Evaluating these multilinear maps in coordinate directions gives the bounds
\[
|\partial^\kappa B_h(x)|
\le d^{\iota_{\mathrm{der}}} C_{B,\iota_{\mathrm{der}}}h^{-\iota_{\mathrm{der}}}
\qquad(|\kappa|=\iota_{\mathrm{der}}\le p_*),
\]
and
\[
|\partial^\kappa B_h(x)-\partial^\kappa B_h(x')|
\le d^{p_*}M_Bh^{-\gamma}\|x-x'\|_2^{s_o}
\qquad(|\kappa|=p_*).
\]
Here the dimension factors are valid because every coordinate direction has norm at most \(d\). Define
\[
S_B=\sum_{\iota_{\mathrm{der}}=0}^{p_*}d^{\iota_{\mathrm{der}}}C_{B,\iota_{\mathrm{der}}},
\qquad
D_B=S_B+d^{p_*}M_B+1.
\]
For \(0<h\le1\) and \(\iota_{\mathrm{der}}\le p_*<\gamma\), one has
\(h^{-\iota_{\mathrm{der}}}\le h^{-\gamma}\). Applying the norm definition in
\cref{obj:def:holder-norm} to the preceding derivative and modulus bounds proves
\[
D_B>0,
\qquad
\|B_h\|_{H^\gamma}\le D_Bh^{-\gamma}
\qquad(0<h\le1).
\]
This also covers the integer smoothness endpoints, where \(s_o=1\).
% lean: CausalSmith.Stat.TwosamplePointcateAnnotationFrontier.oracle_macro_holder

\item Choose the public amplitude constant
\[
\eta_o=\frac{1}{128(D_B+1)}.
\]
It satisfies
\[
\eta_o>0,\qquad
\eta_o(D_B+1)=\frac1{128},\qquad
\eta_o\le\frac18,\qquad
\eta_o\le\frac1{16},\qquad
\eta_oD_B\le L.
\]
For each \(n\ge2\), define
\[
h_n=n^{-1/(2\gamma+d)},
\qquad
b_n=\eta_o r_o(n).
\]
Since \(2\gamma+d>0\) and \(n\ge2\), the power identities give
\[
0<h_n\le1,
\qquad
h_n^\gamma=r_o(n),
\qquad
n\,r_o(n)^2h_n^d=1,
\qquad
0<b_n\le\frac18.
\]
Consequently,
\[
\|b_nB_{h_n}\|_{H^\gamma}
\le b_nD_Bh_n^{-\gamma}
=\eta_oD_B
\le L.
\]

For \(\theta\in\{-1,+1\}\), define the primitive law \(P_{o,\theta}\) by
\[
X\sim\operatorname{Unif}([0,1]^d),
\]
\[
\mathcal L_{P_{o,\theta}}\bigl(A,Y(0),Y(1)\mid X=x\bigr)
=
\operatorname{Bernoulli}(1/2)
\otimes
\operatorname{Bernoulli}(1/2)
\otimes
\operatorname{Bernoulli}\bigl(1/2+\theta b_nB_{h_n}(x)\bigr),
\qquad x\in[0,1]^d.
\]
These conditional margins are Borel. Conditional independence gives exchangeability, and their bounds give
\[
e(x)=\mu_0(x)=\frac12,
\qquad
\frac18\le\frac12+\theta b_nB_{h_n}(x)\le\frac78.
\]
The constant propensity and control functions have Hölder norms at most \(1/2\), while multiplication by either sign preserves the preceding contrast-norm bound. Thus the construction satisfies every condition of \cref{obj:def:primitive-class}.

The designated versions agree with these admissible margins on the cube. Indeed, conditional-margin uniqueness gives almost-sure equality, and continuity together with the full support of the uniform design promotes that equality to pointwise equality. Hence
\[
P_{o,\theta}\in\mathcal M_{1/2},
\qquad
e_{P_{o,\theta}}(x)=\mu_{0,P_{o,\theta}}(x)=\frac12,
\qquad
\tau_{P_{o,\theta}}(x)=\theta b_nB_{h_n}(x)
\quad(x\in[0,1]^d),
\]
so that
\[
\tau_{P_{o,\theta}}(x_0)=\theta b_n.
\]
% lean: CausalSmith.Stat.TwosamplePointcateAnnotationFrontier.oracle_bump_testing_family

\item Define the observed-record and full-experiment laws by
\[
P_{o,\theta}^{\mathrm{obs}}
=\mathcal L_{P_{o,\theta}}(X,A,Y),
\qquad
M_{o,\theta}=\mathsf E_{n,m}(P_{o,\theta}).
\]
Use the common probability reference measure
\[
\nu_o
=
\operatorname{Unif}([0,1]^d)
\otimes\operatorname{Bernoulli}(1/2)
\otimes\operatorname{Bernoulli}(1/2)
\]
on the covariate, treatment, and observed-response coordinates. Relative to this reference, the observed-record density is
\[
\ell_{o,\theta}(x,a_{\mathrm{obs}},y_{\mathrm{obs}})
=
\begin{cases}
1,&a_{\mathrm{obs}}=0,\\
1+2(2y_{\mathrm{obs}}-1)\theta b_nB_{h_n}(x),
&a_{\mathrm{obs}}=1,
\end{cases}
\]
for \(x\in[0,1]^d\) and \(a_{\mathrm{obs}},y_{\mathrm{obs}}\in\{0,1\}\). The amplitude bound implies
\[
\frac12\le\ell_{o,\theta}\le2.
\]
For any two density values at least \(1/2\), the squared sum of their square roots is at least one. The difference-of-squares identity therefore gives, pointwise,
\[
\bigl(\sqrt{\ell_{o,-1}}-\sqrt{\ell_{o,+1}}\bigr)^2
\le
(\ell_{o,-1}-\ell_{o,+1})^2
\le16b_n^2B_{h_n}(x)^2.
\]
The Hellinger normalization in \cref{obj:def:hellinger} can be evaluated under \(\nu_o\): the densities relative to the sum of the two laws are
\(\ell_{o,\theta}/(\ell_{o,-1}+\ell_{o,+1})\), and the common denominator cancels. Thus
\[
\mathsf h^2(P_{o,-1}^{\mathrm{obs}},P_{o,+1}^{\mathrm{obs}})
=
\int
\bigl(\sqrt{\ell_{o,-1}}-\sqrt{\ell_{o,+1}}\bigr)^2\,d\nu_o.
\]
Since \(B_{h_n}^2\le\mathbf1_{C_{h_n}}\), the containment
\(C_{h_n}\subseteq[0,1]^d\) and the product volume of this cube give
\[
\int_{[0,1]^d}B_{h_n}(x)^2\,dx
\le\operatorname{vol}(C_{h_n})
=h_n^d.
\]
Combining the preceding displays proves
\[
\mathsf h^2(P_{o,-1}^{\mathrm{obs}},P_{o,+1}^{\mathrm{obs}})
\le16b_n^2h_n^d.
\]

The auxiliary covariate-treatment marginal is identical under the two laws:
\[
P_{o,-1,XA}=P_{o,+1,XA}
=
\operatorname{Unif}([0,1]^d)\otimes\operatorname{Bernoulli}(1/2).
\]
The product bound and common-factor identity in
\cref{obj:lem:hellinger-block-calculus} therefore yield
\[
\mathsf h^2(M_{o,-1},M_{o,+1})
\le
n\,\mathsf h^2(P_{o,-1}^{\mathrm{obs}},P_{o,+1}^{\mathrm{obs}})
\le16n b_n^2h_n^d
=16\eta_o^2
\le\frac1{16}.
\]
This applies to every \(m\ge0\); at \(m=0\), the auxiliary product is the probability law on the empty record block. The independent uniform randomizer is another common probability factor. Finally, symmetry of the squared square-root difference gives the orientation used below:
\[
\mathsf h^2(M_{o,+1},M_{o,-1})
=
\mathsf h^2(M_{o,-1},M_{o,+1})
\le\frac1{16}.
\]
% lean: CausalSmith.Stat.TwosamplePointcateAnnotationFrontier.oracle_bump_information_at_rate

\item The target is a well-defined functional of these full-experiment laws. To see this, equality of two such laws implies equality of their first labeled-record marginals, since \(n>0\). Under either constructed law, the observed conditional probabilities identify the arm means through
\[
\mu_{1,P_{o,\theta}}(x)
=
2P_{o,\theta}(A=1,Y=1\mid X=x),
\qquad
\mu_{0,P_{o,\theta}}(x)
=
2P_{o,\theta}(A=0,Y=1\mid X=x)
\]
almost everywhere under the uniform design. Thus equality of the observed laws identifies all three conditional margins. Their conditional product construction then identifies the primitive laws, and uniqueness of the canonical continuous contrast gives
\[
M_{o,\theta}=M_{o,\theta'}
\quad\Longrightarrow\quad
\tau_{P_{o,\theta}}(x_0)=\tau_{P_{o,\theta'}}(x_0).
\]
In particular, \(M_{o,+1}\) and \(M_{o,-1}\) are distinct because \(b_n>0\). Define a functional on all measures on the full sample space by
\[
\Psi_o(\mathbb Q_o)=
\begin{cases}
b_n,&\mathbb Q_o=M_{o,+1},\\
-b_n,&\mathbb Q_o=M_{o,-1},\\
0,&\mathbb Q_o\notin\{M_{o,+1},M_{o,-1}\}.
\end{cases}
\]
It satisfies
\[
\Psi_o(M_{o,\theta})=\tau_{P_{o,\theta}}(x_0).
\]

The supplied propensity is the same function for both hypotheses, including its extension outside the cube:
\[
e_o(x)=
\begin{cases}
1/2,&x\in[0,1]^d,\\
0,&x\notin[0,1]^d.
\end{cases}
\]
For any admissible supplied-propensity decision, its section at \(e_o\) is measurable. Apply \cref{obj:lem:published-fuzzy-testing}, the absolute-loss bound of
\citep[Section 3, Lemma 1]{KennedyBalakrishnanRobinsWasserman2024},
with singleton priors, one observation consisting of the entire record block and randomizer, and the choices
\[
\omega_o=\operatorname{Dirac}(*),
\qquad
\mathbb B_*^{\,o}=M_{o,+1},
\qquad
\mathbb C_*^{\,o}=M_{o,-1},
\qquad
\mathcal P_o^{\mathrm{block}}=\{M_{o,+1},M_{o,-1}\},
\]
\[
v_{\mathrm{test}}=1,
\qquad
s_{\mathrm{test}}=2b_n,
\qquad
u_{\mathrm{test}}=\frac1{16}.
\]
The separation is \(2b_n>0\), and the preceding Hellinger bound supplies the distance hypothesis. Consequently,
\[
\sup_{\theta\in\{-1,+1\}}
\mathcal R^e_{n,m}(T^e,P_{o,\theta})
\ge
\frac{b_n}{2}
\left(1-\frac{\sqrt{63}}{32}\right).
\]
Since \(63\le64\),
\[
1-\frac{\sqrt{63}}{32}\ge\frac34,
\qquad
\frac{b_n}{2}
\left(1-\frac{\sqrt{63}}{32}\right)\ge\frac38b_n.
\]
The two primitive laws belong to \(\mathcal M_{1/2}\), so taking the infimum over decisions proves
\[
c=\frac38\eta_o>0,
\qquad
c\,r_o(n)\le R^e_{1/2}(n,m)
\qquad(n\ge2,\ m\ge0).
\]
% lean: CausalSmith.Stat.TwosamplePointcateAnnotationFrontier.oracle_subfamily_lower_rate

\item We next establish the upper bound. For every \(h>0\), define the localization measure and weight by
\[
d\nu_h(x)=h^{-d}\mathbf1_{C_h}(x)\,dx,
\qquad
w_h(x)=h^{-d}\mathbf1_{C_h}(x).
\]
The measure \(\nu_h\) is a probability measure because
\(\operatorname{vol}(C_h)=h^d\). For \(0<h\le1\), its support lies in the design cube.

To make the coarse basis explicit, define
\[
\mathcal I_2=
\{\kappa\in\mathbb N^d:|\kappa|\le2\},
\qquad
q=|\mathcal I_2|,
\]
and the one-dimensional polynomials
\[
\mathsf P_0(\upsilon)=1,\qquad
\mathsf P_1(\upsilon)=\sqrt{12}\,\upsilon,\qquad
\mathsf P_2(\upsilon)=\sqrt5\,(6\upsilon^2-1/2),
\qquad \upsilon\in\mathbb R.
\]
At the current localization width \(h\), the basis and evaluation vector are
\[
r_\kappa(x)=
\prod_{\iota_{\mathrm{cov}}=1}^d
\mathsf P_{\kappa_{\iota_{\mathrm{cov}}}}\bigl((x_{\iota_{\mathrm{cov}}}-1/2)/h\bigr),
\qquad
r(x)=(r_\kappa(x))_{\kappa\in\mathcal I_2},
\qquad
r_0=r(x_0),
\]
for every \(x\in\mathbb R^d\). Define the evaluation kernel and its dimension constant by
\[
\mathcal K_h(x)=r_0^\top r(x),
\qquad
K_{\mathrm{ev}}=q(4^d)^2.
\]
The constant coordinate belongs to \(\mathcal I_2\), so \(q\ge1\) and \(K_{\mathrm{ev}}>0\).

Direct integration of the three one-dimensional formulas gives
\[
\int_{-1/2}^{1/2}
\mathsf P_{j_{\mathrm{pol}}}(\upsilon)
\mathsf P_{j'_{\mathrm{pol}}}(\upsilon)\,d\upsilon
=
\mathbf1_{\{j_{\mathrm{pol}}=j'_{\mathrm{pol}}\}}
\qquad
(j_{\mathrm{pol}},j'_{\mathrm{pol}}\in\{0,1,2\}).
\]
Scaling each coordinate and multiplying these identities yields
\[
\int r_\kappa(x)r_{\kappa'}(x)\,d\nu_h(x)
=
\mathbf1_{\{\kappa=\kappa'\}}.
\]
Expanding the finite sums consequently gives polynomial reproduction and kernel energy:
\[
\int\mathcal K_h(x)\,r(x)^\top v\,d\nu_h(x)
=r_0^\top v
\qquad(v\in\mathbb R^q),
\]
\[
\int\mathcal K_h(x)^2\,d\nu_h(x)=\|r_0\|_2^2.
\]
The displayed polynomial formulas also give
\(|\mathsf P_{j_{\mathrm{pol}}}(\upsilon)|\le4\) on \([-1/2,1/2]\). Thus
\[
|r_\kappa(x)|\le4^d\quad(x\in C_h),
\qquad
|\mathcal K_h(x)|\le K_{\mathrm{ev}}\quad(x\in C_h),
\qquad
\|r_0\|_2^2\le K_{\mathrm{ev}}.
\]
All these identities hold for every \(h>0\).
% lean: CausalSmith.Stat.TwosamplePointcateAnnotationFrontier.oracle_polynomial_reproduction

\item Fix \(P\in\mathcal M\). Define the supplied-propensity response, its localized record contribution, the raw average, and its mean by
\[
Z_P(x,a_{\mathrm{obs}},y_{\mathrm{obs}})
=
\frac{a_{\mathrm{obs}}y_{\mathrm{obs}}}{e_P(x)}
-
\frac{(1-a_{\mathrm{obs}})y_{\mathrm{obs}}}{1-e_P(x)},
\]
\[
F_{P,h}(x,a_{\mathrm{obs}},y_{\mathrm{obs}})
=
w_h(x)\mathcal K_h(x)
Z_P(x,a_{\mathrm{obs}},y_{\mathrm{obs}}),
\]
\[
\overline T_{P,h}(\mathcal D_{n,m},U)
=
\frac1n\sum_{i=1}^nF_{P,h}(X_i,A_i,Y_i),
\qquad
\mu_{P,h}
=
\int\overline T_{P,h}\,d\mathsf E_{n,m}(P).
\]
These definitions apply to all \(x\in\mathbb R^d\), binary record coordinates, and every dataset and randomizer, using the quotient convention stated above.

On the design cube, overlap and the binary record coordinates give
\[
|Z_P(x,a_{\mathrm{obs}},y_{\mathrm{obs}})|\le\varepsilon^{-1}.
\]
Outside the cube, the zero extension of the propensity gives
\[
Z_P(x,a_{\mathrm{obs}},y_{\mathrm{obs}})
=-(1-a_{\mathrm{obs}})y_{\mathrm{obs}},
\qquad
|Z_P(x,a_{\mathrm{obs}},y_{\mathrm{obs}})|\le1\le\varepsilon^{-1}.
\]
Hence the response bound is global, and
\[
|F_{P,h}(x,a_{\mathrm{obs}},y_{\mathrm{obs}})|
\le\varepsilon^{-1}h^{-d}K_{\mathrm{ev}}.
\]
Borel measurability and this bound supply all integrability and finite second moments used below.

For \(0<h\le1\), exchangeability and the conditional-margin identities in
\cref{obj:def:primitive-class}, followed by cancellation of the overlap denominators, give
\[
\mathbb E_P\left[
w_h(X)\mathcal K_h(X)\frac{AY}{e_P(X)}
\right]
=
\int_{[0,1]^d}w_h(x)\mathcal K_h(x)\mu_{1,P}(x)\,dx,
\]
\[
\mathbb E_P\left[
w_h(X)\mathcal K_h(X)\frac{(1-A)Y}{1-e_P(X)}
\right]
=
\int_{[0,1]^d}w_h(x)\mathcal K_h(x)\mu_{0,P}(x)\,dx.
\]
Subtracting and using the canonical contrast identity gives
\[
\mathbb E_P[F_{P,h}(X,A,Y)]
=
\int\mathcal K_h(x)\tau_P(x)\,d\nu_h(x).
\]
Every labeled coordinate has this expectation, while the auxiliary block and randomizer integrate out as probability factors. Since \(n>0\),
\[
\mu_{P,h}
=
\int\mathcal K_h(x)\tau_P(x)\,d\nu_h(x).
\]
% lean: CausalSmith.Stat.TwosamplePointcateAnnotationFrontier.oracle_average_mean

\item We now bound this mean's bias over the full range \(0<h\le1\). Define the Taylor polynomial
\[
p(x)=
\sum_{|\kappa|\le p_*}
\frac{\partial^\kappa\tau_P(x_0)}{\kappa!}
(x-x_0)^\kappa,
\qquad
\kappa!=\prod_{\iota_{\mathrm{cov}}=1}^d\kappa_{\iota_{\mathrm{cov}}}!.
\]
Because \(p_*\le2\), the coarse polynomial basis spans \(p\). Thus there exists a coefficient vector \(\vartheta_{\mathrm{Tay}}\in\mathbb R^q\) with
\[
p_{\vartheta_{\mathrm{Tay}}}(x)
=
r(x)^\top\vartheta_{\mathrm{Tay}}
=
p(x)
\qquad(x\in\mathbb R^d),
\qquad
p_{\vartheta_{\mathrm{Tay}}}(x_0)=\tau_P(x_0).
\]

If \(\gamma=1\), then \(p_*=0\) and \(p(x)=\tau_P(x_0)\). The Hölder norm condition gives the Lipschitz bound, and the geometry of \(C_h\) gives
\[
|\tau_P(x)-p(x)|
\le L\|x-x_0\|_2
\le L\sqrt d\,h
\le L(\sqrt d+1)h
\qquad(x\in C_h).
\]

If \(1<\gamma\le3\), then \(1\le p_*\le2\). For each \(x\in C_h\), define the segment function
\[
g_{P,x}(\sigma_{\mathrm{seg}})=\tau_P\bigl(x_0+\sigma_{\mathrm{seg}}(x-x_0)\bigr),
\qquad 0\le\sigma_{\mathrm{seg}}\le1.
\]
The segment lies in the design cube. Expanding its \(p_*\)-th derivative in coordinate partials and using their \(s_o\)-Hölder modulus gives
\[
\left|g_{P,x}^{(p_*)}(\sigma_{\mathrm{seg}})-g_{P,x}^{(p_*)}(0)\right|
\le
(d+1)^{p_*}L
\|x-x_0\|_2^{p_*+s_o}\sigma_{\mathrm{seg}}^{s_o}.
\]
Indeed, the multinomial expansion contributes at most
\((d+1)^{p_*}\|x-x_0\|_2^{p_*}\), and each top-partial difference is bounded by
\(L\sigma_{\mathrm{seg}}^{s_o}\|x-x_0\|_2^{s_o}\).
The one-dimensional Lagrange remainder at degree \(p_*-1\), followed by subtraction of the degree-\(p_*\) Taylor term, supplies some
\(\xi_{P,x}\in(0,1)\) such that
\[
\tau_P(x)-p(x)
=
\frac{
g_{P,x}^{(p_*)}(\xi_{P,x})
-g_{P,x}^{(p_*)}(0)
}{p_*!}.
\]
Since \(\xi_{P,x}^{s_o}\le1\) and \(p_*!\ge1\),
\[
|\tau_P(x)-p(x)|
\le(d+1)^{p_*}L\|x-x_0\|_2^\gamma.
\]
Using \(\|x-x_0\|_2\le(d+1)h\) and
\(\gamma=p_*+s_o\le p_*+1\) yields
\[
|\tau_P(x)-p(x)|
\le L(d+1)^{2p_*+1}h^\gamma
\qquad(x\in C_h).
\]

Define the positive public constants
\[
C_{\mathrm{Tay}}=
\begin{cases}
L(\sqrt d+1),&\gamma=1,\\
L(d+1)^{2p_*+1},&1<\gamma\le3,
\end{cases}
\qquad
C_{\mathrm{bias}}=K_{\mathrm{ev}}C_{\mathrm{Tay}}.
\]
Both cases establish
\[
\sup_{x\in C_h}|\tau_P(x)-p_{\vartheta_{\mathrm{Tay}}}(x)|
\le C_{\mathrm{Tay}}h^\gamma.
\]
Polynomial reproduction and the mean identity therefore imply
\[
\begin{aligned}
|\mu_{P,h}-\tau_P(x_0)|
&=
\left|
\int\mathcal K_h(x)
\bigl(\tau_P(x)-p_{\vartheta_{\mathrm{Tay}}}(x)\bigr)
\,d\nu_h(x)
\right|\\
&\le
\int|\mathcal K_h(x)|
\left|\tau_P(x)-p_{\vartheta_{\mathrm{Tay}}}(x)\right|
\,d\nu_h(x)\\
&\le C_{\mathrm{bias}}h^\gamma.
\end{aligned}
\]
At \(h=h_n\), this becomes
\[
|\mu_{P,h_n}-\tau_P(x_0)|
\le C_{\mathrm{bias}}r_o(n).
\]
% lean: CausalSmith.Stat.TwosamplePointcateAnnotationFrontier.oracle_tuned_average_bias

\item The global response bound gives the recordwise squared bound
\[
F_{P,h}(X,A,Y)^2
\le\varepsilon^{-2}\bigl(w_h(X)\mathcal K_h(X)\bigr)^2.
\]
Under the uniform design, localization and the kernel-energy identity give
\[
\begin{aligned}
\mathbb E_P[F_{P,h}(X,A,Y)^2]
&\le
\varepsilon^{-2}\int_{[0,1]^d}
\bigl(w_h(x)\mathcal K_h(x)\bigr)^2\,dx\\
&=
\varepsilon^{-2}h^{-d}\int\mathcal K_h(x)^2\,d\nu_h(x)\\
&=
\varepsilon^{-2}h^{-d}\|r_0\|_2^2.
\end{aligned}
\]
Independence of the \(n\) labeled records yields
\[
\begin{aligned}
\int(\overline T_{P,h}-\mu_{P,h})^2\,d\mathsf E_{n,m}(P)
&=
\frac1n\operatorname{Var}_P(F_{P,h}(X,A,Y))\\
&\le\frac1n\mathbb E_P[F_{P,h}(X,A,Y)^2]\\
&\le\frac{\varepsilon^{-2}\|r_0\|_2^2}{nh^d}.
\end{aligned}
\]
The first inequality uses
\[
\operatorname{Var}_P(F_{P,h})
=
\mathbb E_P[F_{P,h}^2]
-\bigl(\mathbb E_P[F_{P,h}]\bigr)^2
\le\mathbb E_P[F_{P,h}^2].
\]
Cauchy--Schwarz on the probability experiment then gives
\[
\int|\overline T_{P,h}-\mu_{P,h}|\,d\mathsf E_{n,m}(P)
\le
\left(
\int(\overline T_{P,h}-\mu_{P,h})^2\,d\mathsf E_{n,m}(P)
\right)^{1/2}
\le
\sqrt{\frac{\varepsilon^{-2}\|r_0\|_2^2}{nh^d}}.
\]
At \(h=h_n\), use the target-coordinate bound and
\[
(nh_n^d)^{-1/2}=r_o(n)
\]
to obtain
\[
\int|\overline T_{P,h_n}-\mu_{P,h_n}|\,d\mathsf E_{n,m}(P)
\le\sqrt{\varepsilon^{-2}K_{\mathrm{ev}}}\,r_o(n)
\le C_{\mathrm{dev}}r_o(n),
\qquad
C_{\mathrm{dev}}=\sqrt{\varepsilon^{-2}K_{\mathrm{ev}}}+1>0.
\]
The estimate is uniform in \(P\), \(n\ge2\), and \(m\ge0\).
% lean: CausalSmith.Stat.TwosamplePointcateAnnotationFrontier.oracle_tuned_absolute_deviation

\item The probability-valued arm-mean conditions in \cref{obj:def:primitive-class} imply
\[
-1
\le
\mu_{1,P}(x_0)-\mu_{0,P}(x_0)
=
\tau_P(x_0)
\le1.
\]
The displayed smoother equals
\[
T^e(\mathcal D_{n,m},U,e_P)
=
\operatorname{clip}_{[-1,1]}
\bigl(\overline T_{P,h_n}(\mathcal D_{n,m},U)\bigr).
\]
Projection onto this interval contracts distance to any target inside it. Thus, for every dataset and randomizer,
\[
\begin{aligned}
|T^e(\mathcal D_{n,m},U,e_P)-\tau_P(x_0)|
&\le|\overline T_{P,h_n}(\mathcal D_{n,m},U)-\tau_P(x_0)|\\
&\le
|\overline T_{P,h_n}(\mathcal D_{n,m},U)-\mu_{P,h_n}|
+
|\mu_{P,h_n}-\tau_P(x_0)|.
\end{aligned}
\]
Integrating and applying the preceding two bounds gives
\[
\mathcal R^e_{n,m}(T^e,P)
\le
(C_{\mathrm{dev}}+C_{\mathrm{bias}})r_o(n).
\]
Define
\[
C_{\mathrm{up}}=C_{\mathrm{dev}}+C_{\mathrm{bias}}>0.
\]
The bounded measurable raw average and continuous clipping map give measurability of the supplied-propensity section. The clipping formula itself gives the pointwise guarantee
\[
\left|T^e(\mathcal D_{n,m},U,e_P)\right|\le1
\]
for every dataset and randomizer.
% lean: CausalSmith.Stat.TwosamplePointcateAnnotationFrontier.oracle_smoother_upper_rate

\item For completeness, the smoother has a total admissible extension to arbitrary supplied functions. Define
\[
\widetilde T^e(\mathcal D_{n,m},U,e)=
\begin{cases}
T^e(\mathcal D_{n,m},U,e),&e\text{ is Borel},\\
0,&e\text{ is not Borel}.
\end{cases}
\]
For every fixed Borel \(e\), the totalized reciprocal, record-coordinate maps, finite sum, and clipping are Borel; the other section is constant. Since each designated \(e_P\) is Borel, this extension agrees with the displayed smoother at every admitted law. It is therefore a candidate in \cref{obj:def:oracle-risk}, and the uniform bound proves
\[
R^e(n,m)
\le
\sup_{P\in\mathcal M}
\mathcal R^e_{n,m}(\widetilde T^e,P)
\le C_{\mathrm{up}}r_o(n).
\]

Finally, retain the lower constant already constructed and set
\[
C=\max\{c,C_{\mathrm{up}}\}.
\]
These are finite real constants depending on the fixed public parameters, with
\[
0<c\le C.
\]
Because \(r_o(n)\ge0\), enlarging the upper constant preserves both the minimax and individual-decision bounds. Combining the lower bound with subfamily monotonicity gives, for every \(n\ge2\) and \(m\ge0\),
\[
c\,r_o(n)
\le R^e_{1/2}(n,m)
\le R^e(n,m)
\le C\,r_o(n),
\]
and, for every \(P\in\mathcal M\),
\[
\mathcal R^e_{n,m}(T^e,P)\le C\,r_o(n).
\]
Together with the measurability and pointwise clipping bound above, these are all the asserted conclusions.
% lean: CausalSmith.Stat.TwosamplePointcateAnnotationFrontier.supplied_propensity
\end{enumerate}
\end{proof}
